\section*{Problem 1}

一个人工神经元的输入为
\[
\bm{x} = (1, -2, 3)^\top,\quad
\bm{w} = (2, -1, 0.5)^\top,\quad
b = -1.
\]
记净输入为 $z = \bm{w}^\top \bm{x} + b$。
\begin{enumerate}
    \item 计算该神经元的净输入 $z$。
    \item 若激活函数为阶梯函数
    \[
    h(z) = \begin{cases}
        1, & z \ge 0,\\
        0, & z < 0,
    \end{cases}
    \]
    求神经元输出。
    \item 若激活函数分别为 $\operatorname{ReLU}(z) = \max\{0, z\}$ 和 LeakyReLU，其中
    \[
    \operatorname{LeakyReLU}(z) = \begin{cases}
        z,      & z \ge 0,\\
        0.01z,  & z < 0,
    \end{cases}
    \]
    分别求神经元输出，并说明二者在 $z < 0$ 时的主要区别。
\end{enumerate}

\begin{proof}
    \begin{enumerate}
        \item 净输入为
        \[
            z = 2\cdot 1 + (-1)(-2) + 0.5\cdot 3 - 1
              = \boxed{4.5}.
        \]
        \item 因为 $z=4.5>0$，阶梯函数的输出为
        \[
            \boxed{h(z)=1}.
        \]
        \item 此处净输入为正，故两种激活函数均输出
        \[
            \operatorname{ReLU}(4.5)
            = \operatorname{LeakyReLU}(4.5)
            = \boxed{4.5}.
        \]
        当 $z<0$ 时，ReLU 的输出恒为 $0$，而 LeakyReLU 保留 $0.01z$，因而仍保留一条斜率为 $0.01$ 的负半轴分支。
    \end{enumerate}
\end{proof}

\section*{Problem 2}

Logistic 激活函数定义为
\[
\sigma(z) = \frac{1}{1 + \exp(-z)}.
\]
已知其导数为 $\sigma'(z) = \sigma(z)(1 - \sigma(z))$。
\begin{enumerate}
    \item 计算 $\sigma(0)$ 与 $\sigma'(0)$。
    \item 近似计算 $\sigma(4), \sigma'(4), \sigma(-4), \sigma'(-4)$。可使用 $\exp(-4) \approx 0.0183$, $\exp(4) \approx 54.6$。
    \item 根据以上结果，解释为什么 Logistic 激活函数在 $|z|$ 较大时容易导致梯度消失。
\end{enumerate}

\begin{proof}
    \begin{enumerate}
        \item 代入 $z=0$，得到
        \[
            \sigma(0)=\frac12,\qquad
            \sigma'(0)=\frac12\left(1-\frac12\right)=\boxed{\frac14}.
        \]
        \item 利用题给近似值，
        \[
        \begin{aligned}
            \sigma(4)
            &\approx \frac{1}{1+0.0183}
             \approx 0.9820, \\
            \sigma'(4)
            &\approx 0.9820(1-0.9820)
             \approx 0.0177, \\
            \sigma(-4)
            &\approx \frac{1}{1+54.6}
             \approx 0.0180, \\
            \sigma'(-4)
            &\approx 0.0180(1-0.0180)
             \approx 0.0177.
        \end{aligned}
        \]
        \item 当 $z\to+\infty$ 时，$\sigma(z)\to1$；当 $z\to-\infty$ 时，$\sigma(z)\to0$。两种情形下
        \[
            \sigma'(z)=\sigma(z)(1-\sigma(z))\to0.
        \]
        反向传播中的梯度需要逐层乘以该导数，多层连乘后会迅速趋近于零，因此前层参数难以获得有效更新。
    \end{enumerate}
\end{proof}

\section*{Problem 3}

考虑一个使用 ReLU 激活函数的神经元
\[
z = \bm{w}^\top \bm{x} + b,\quad a = \operatorname{ReLU}(z),
\]
设训练集中所有样本都满足 $\bm{x} \ge 0$（逐分量非负）。
\begin{enumerate}
    \item 写出 ReLU 的函数表达式及其在 $z \neq 0$ 处的导数。
    \item 设 $\bm{w} = (-1, -2)^\top$, $b = 0$，对任意满足 $\bm{x} \ge 0$ 的样本，说明为什么有 $z \le 0$。
    \item 若某个神经元对所有训练样本都有 $z < 0$，说明其权重为什么难以通过梯度下降更新，并说明 LeakyReLU 为什么可以缓解这一问题。
\end{enumerate}

\begin{proof}
    \begin{enumerate}
        \item ReLU 及其在 $z\neq0$ 处的导数分别为
        \[
            \operatorname{ReLU}(z)=
            \begin{cases}
                z, & z>0,\\
                0, & z\le0,
            \end{cases}
            \qquad
            \operatorname{ReLU}'(z)=
            \begin{cases}
                1, & z>0,\\
                0, & z<0.
            \end{cases}
        \]
        \item 写成 $\bm{x}=(x_1,x_2)^\top$。由 $x_1,x_2\ge0$，有
        \[
            z=-x_1-2x_2\le0.
        \]
        \item 若所有训练样本均满足 $z<0$，则 $\operatorname{ReLU}'(z)=0$。以损失函数 $L$ 为例，链式法则给出
        \[
            \frac{\partial L}{\partial \bm{w}}
            = \frac{\partial L}{\partial a}
              \operatorname{ReLU}'(z)\bm{x}
            = \bm{0},
            \qquad
            \frac{\partial L}{\partial b}=0.
        \]
        因而梯度下降无法把该神经元带回激活区。LeakyReLU 在 $z<0$ 时的导数为 $0.01$，梯度虽小但不为零，所以参数仍可能得到更新。
    \end{enumerate}
\end{proof}
